\documentclass[a4paper,11pt]{article}

% ── Packages ──
\usepackage[utf8]{inputenc}
\usepackage[T1]{fontenc}
\usepackage[french]{babel}
\usepackage{geometry}
\geometry{margin=2.2cm, top=2.5cm, bottom=2.5cm}
\usepackage{xcolor}
\usepackage{listings}
\usepackage{fancyhdr}
\usepackage{titlesec}
\usepackage{hyperref}
\usepackage{tocloft}
\usepackage{amsmath}
\usepackage{graphicx}
\usepackage{tabularx}
\usepackage{booktabs}
\usepackage{colortbl}
\usepackage{enumitem}
\usepackage{fancybox}
\usepackage{mdframed}

% ── Couleurs ──
\definecolor{codeblue}{HTML}{1F4E79}
\definecolor{codebg}{HTML}{F5F7FA}
\definecolor{codeborder}{HTML}{C5D0DC}
\definecolor{codecomment}{HTML}{2E8B57}
\definecolor{codestring}{HTML}{A31515}
\definecolor{codekeyword}{HTML}{0000FF}
\definecolor{sectionblue}{HTML}{1F4E79}
\definecolor{lightblue}{HTML}{D6E4F0}
\definecolor{darkgray}{HTML}{444444}
\definecolor{tabheader}{HTML}{1F4E79}

% ── Listings: style production ──
\lstdefinestyle{cstyle}{
  language=C,
  backgroundcolor=\color{codebg},
  basicstyle=\ttfamily\small,
  keywordstyle=\color{codekeyword}\bfseries,
  commentstyle=\color{codecomment}\itshape,
  stringstyle=\color{codestring},
  numberstyle=\tiny\color{gray},
  numbers=left,
  numbersep=8pt,
  frame=single,
  rulecolor=\color{codeborder},
  breaklines=true,
  breakatwhitespace=false,
  tabsize=2,
  showstringspaces=false,
  captionpos=b,
  xleftmargin=1.2em,
  framexleftmargin=1.2em,
  aboveskip=10pt,
  belowskip=6pt,
  morekeywords={float32_t,uint16_t,uint32_t,uint8_t,int32_t,arm_rfft_fast_instance_f32,
    ADC_HandleTypeDef,DMA_HandleTypeDef,DAC_HandleTypeDef,TIM_HandleTypeDef,
    UART_HandleTypeDef,GPIO_InitTypeDef,ADC_ChannelConfTypeDef,DAC_ChannelConfTypeDef,
    volatile},
  literate={é}{{\'e}}1 {è}{{\`e}}1 {ê}{{\^e}}1 {à}{{\`a}}1 {â}{{\^a}}1
            {î}{{\^i}}1 {ô}{{\^o}}1 {ù}{{\`u}}1 {û}{{\^u}}1 {ç}{{\c{c}}}1
            {→}{{$\rightarrow$}}1,
}

\lstset{style=cstyle}

% ── Headers/footers ──
\pagestyle{fancy}
\fancyhf{}
\fancyhead[L]{\small\textit{\color{darkgray}Radar Doppler 24\,GHz -- Documentation firmware}}
\fancyhead[R]{\small\textit{\color{darkgray}Projet L3 EEA}}
\fancyfoot[C]{\small\color{darkgray}\thepage}
\renewcommand{\headrulewidth}{0.4pt}
\renewcommand{\footrulewidth}{0pt}

% ── Titres ──
\titleformat{\section}
  {\Large\bfseries\color{sectionblue}}
  {\thesection}{0.8em}{}
  [\vspace{-0.5em}\textcolor{sectionblue}{\hrule height 0.8pt}]

\titleformat{\subsection}
  {\large\bfseries\color{sectionblue}}
  {\thesubsection}{0.7em}{}

\titleformat{\subsubsection}
  {\normalsize\bfseries\color{sectionblue}}
  {\thesubsubsection}{0.6em}{}

% ── Hyperlinks ──
\hypersetup{
  colorlinks=true,
  linkcolor=sectionblue,
  urlcolor=sectionblue,
  pdfauthor={Yanis},
  pdftitle={Radar Doppler 24 GHz - Documentation firmware},
}

% ── Commande raccourci ──
\newcommand{\code}[1]{\texttt{\small #1}}

\begin{document}

% ══════════════════════════════════════════
% PAGE DE TITRE
% ══════════════════════════════════════════
\begin{titlepage}
\centering
\vspace*{4cm}
{\Huge\bfseries\color{sectionblue} Radar Doppler 24\,GHz\par}
\vspace{0.8cm}
{\Large Documentation du firmware STM32F767ZI\par}
\vspace{0.4cm}
{\large Explication section par section de \texttt{main.c}\par}
\vspace{2cm}
{\large\bfseries Projet L3 EEA\par}
\vspace{0.5cm}
{\large Mars 2026\par}
\vfill
\textcolor{darkgray}{\small Ce document couvre l'ensemble du code applicatif du radar, avec le code source
comment\'e et les explications des concepts sous-jacents (FFT, DMA, fenetrage, Doppler).}
\end{titlepage}

% ══════════════════════════════════════════
% TABLE DES MATIERES
% ══════════════════════════════════════════
\tableofcontents
\newpage

% ══════════════════════════════════════════
% CONVENTIONS
% ══════════════════════════════════════════
\section{Conventions de ce document}

Les extraits de code de \textbf{production} (n\'ecessaires au fonctionnement du radar) sont pr\'esent\'es sur fond bleu clair :

\begin{lstlisting}[style=cstyle,numbers=none]
// Exemple de code production
arm_rfft_fast_f32(&fft_instance, fft_in, fft_out, 0);
\end{lstlisting}

Les extraits de code li\'es au DAC sont d\'esormais document\'es comme partie int\'egrante du syst\`eme :

\begin{lstlisting}[style=cstyle,numbers=none]
// Exemple de code DAC
update_dac_sweep();
\end{lstlisting}

\newpage

% ══════════════════════════════════════════
% 1. VUE D'ENSEMBLE
% ══════════════════════════════════════════
\section{Vue d'ensemble de l'architecture}

Le firmware r\'ealise un v\'elocim\`etre Doppler en quatre \'etapes :

\begin{enumerate}[leftmargin=2em]
  \item \textbf{Acquisition} : l'ADC1 \'echantillonne le signal radar \`a \textbf{100\,kHz}, pilot\'e par TIM2. Le DMA remplit un buffer circulaire de \textbf{2048} mots (interruptions half/full tous les 1024 points).
  \item \textbf{Accumulation + streaming} : chaque demi-buffer (1024 points) est copi\'e vers \texttt{fft\_in} et envoy\'e en \textbf{hop UDP} fragment\'e vers le PC (port 5555).
  \item \textbf{Traitement FFT} : suppression DC, fen\^etrage de Hann, FFT r\'eelle \textbf{4096 points}, recherche de pic sur \texttt{[MIN\_BIN..MAX\_BIN]}, estimation bruit/SNR et conversion Doppler en vitesse.
  \item \textbf{Sorties} : trame binaire UART3 (compatibilit\'e legacy), mise \`a jour OLED (5 Hz), LEDs d'\'etat, et surveillance p\'eriodique du lien Ethernet.
\end{enumerate}

\subsection{Tableau r\'ecapitulatif des p\'eriph\'eriques}

\begin{center}
\small
\begin{tabularx}{\textwidth}{l l X}
\toprule
\rowcolor{lightblue}
\textbf{P\'eriph\'erique} & \textbf{R\^ole} & \textbf{Configuration} \\
\midrule
ADC1           & \'Echantillonnage signal radar     & PA0, 12 bits, d\'eclench\'e par TIM2 TRGO \\
DMA2 Stream0   & Transfert ADC $\to$ m\'emoire         & Mode circulaire, half/full callbacks \\
TIM2           & Horloge ADC (runtime)              & Prescaler=0, ARR=1079 $\rightarrow$ 100\,kHz (reconfig en USER CODE) \\
UART3          & Sortie r\'esultats / debug          & 921600~bauds, trames binaires compatibilit\'e legacy \\
\rowcolor{lightblue}
Ethernet RMII  & Streaming ADC vers PC               & UDP brut, hops de 1024 \'echantillons, port 5555 \\
\rowcolor{lightblue}
OLED SSD1306   & Affichage local                     & I2C1, rafra\^ichissement 200 ms \\
DAC1           & G\'en\'eration signal de r\'ef\'erence    & PA4, DMA1 Stream5 Ch7, TIM4 100\,kHz \\
TIM4           & Horloge DAC (runtime)                 & Prescaler=0, ARR=1079 $\rightarrow$ 100\,kHz \\
\bottomrule
\end{tabularx}

\subsection{Mises \`a jour globales du projet (STM32 + PC)}

Depuis la version initiale, le syst\`eme a \'evolu\'e de mani\`ere significative c\^ot\'e acquisition, traitement et interface :

\begin{itemize}[leftmargin=2em]
  \item \textbf{Cha\^ine STM32 $\rightarrow$ PC modernis\'ee} :
  le firmware diffuse les donn\'ees ADC vers le PC en \textbf{hops UDP} (fen\^etre glissante), ce qui permet un traitement plus fr\'equent et une meilleure r\'eactivit\'e qu'un envoi trame compl\`ete unique.

  \item \textbf{Analyseur PC \texttt{pc\_analyzer} acc\'el\'er\'e} :
  FFT glissante avec fr\'equence de rafra\^ichissement plus \'elev\'ee, axe fr\'equence lisible, curseur de lecture, affichage de vitesse brute et vitesse filtr\'ee.

  \item \textbf{D\'etection et suivi renforc\'es} :
  CFAR r\'eglable en live, lissage spectral, stabilisation du pic, filtre de Kalman configurable, suivi multi-pics (\texttt{Top-N tracks}) et m\'etriques de confiance (stabilit\'e SNR, \^age de piste, jitter, ratio de pertes).

  \item \textbf{Gestion du bruit et perturbations} :
  filtrage secteur (50~Hz et harmoniques), plancher de bruit adaptatif, notch configurables (centre/largeur/profondeur), options de calibration automatique (offset DC, normalisation de gain, profil de bruit idle).

  \item \textbf{Fonctions d'exploitation} :
  mode \textbf{record/replay} pour rejouer des captures hors ligne et comparer des r\'eglages A/B.

  \item \textbf{Interface utilisateur remani\'ee} :
  menu sup\'erieur \texttt{Live Controls} + \texttt{Help}, organisation des r\'eglages par contexte (acquisition, CFAR, bruit/notch, tracking, calibration), aides int\'egr\'ees.

  \item \textbf{Modes de layout} :
  vues focalis\'ees (\texttt{FFT only}, \texttt{Scope only}, \texttt{Spectrogram only}) et ratios de panneaux r\'eglables.

  \item \textbf{Scope long + \'echelle verticale par division} :
  \texttt{Scope ms/div} monte jusqu'\`a \textbf{1000 ms/div} avec historique 1 s (\textbf{100000 \'echantillons} \`a 100~kHz), et r\'eglage \texttt{Scope amplitude/div (ADC)} en mode manuel (\(\pm 4\) divisions).

  \item \textbf{Spectrogramme type waterfall} :
  affichage vertical temporel stabilis\'e (ancien en haut, r\'ecent en bas) pour une lecture plus naturelle de l'\'evolution Doppler.

  \item \textbf{Portage \texttt{pc\_scope}} :
  une partie des am\'eliorations UI/traitement a \'et\'e report\'ee dans \texttt{pc\_scope} (menus, layout, bruit/notch, rendu spectrogramme), avec poursuite de l'alignement fonctionnel avec \texttt{pc\_analyzer}.
\end{itemize}

Note : cette documentation reste centr\'ee sur \texttt{main.c}. La liste ci-dessus synth\'etise les \'evolutions syst\`eme introduites c\^ot\'e applications PC pour l'exploitation radar temps r\'eel.
\end{center}

\subsection{Sch\'ema de flux des donn\'ees}

\begin{center}
\fbox{\parbox{0.9\textwidth}{\centering\small
\texttt{PA0 (signal radar)} $\xrightarrow{\text{ADC1 @ 100kHz}}$ \texttt{adc\_buf[2048]} $\xrightarrow{\text{DMA half/full (1024)}}$ \texttt{fft\_in[4096]} $\xrightarrow{\text{FFT}}$ \texttt{fft\_mag[2048]} $\xrightarrow{\text{pic/SNR}}$ \texttt{UDP + UART + OLED}
}}
\end{center}

\newpage

% ══════════════════════════════════════════
% 2. INCLUDES ET CONSTANTES
% ══════════════════════════════════════════
\section{Includes et constantes}

\subsection{Fichiers inclus}

\begin{lstlisting}[style=cstyle,caption={Fichiers d'en-t\^ete inclus}]
#include "main.h"       // Genere par CubeMX : defines des pins,
                        // prototypes HAL, config peripheriques
#include "string.h"     // memcpy, memset (standard C)
#include <math.h>       // powf, sinf, fmodf, log10f
                        // pour le calcul FFT et sweep DAC
#include <stdio.h>      // sprintf pour formater les messages UART
#include "arm_math.h"   // Bibliotheque CMSIS-DSP d'ARM :
                        // fournit les fonctions FFT, multiplication
                        // vectorielle, moyenne, max, cos
                        // optimisees pour Cortex-M7 (SIMD, FPU)
\end{lstlisting}

\textbf{Pourquoi CMSIS-DSP ?} La biblioth\`eque CMSIS-DSP exploite le FPU simple pr\'ecision et les instructions SIMD du Cortex-M7 pour acc\'el\'erer les calculs vectoriels. Dans ce projet, elle permet une FFT r\'eelle 4096 points + post-traitement dans la contrainte temps r\'eel (\(\sim\) 24 trames/s).

\subsection{Constantes de production}

\begin{lstlisting}[style=cstyle,caption={D\'efinitions des constantes de production (align\'ees avec main.c)}]
#define ADC_BUF_SIZE      2048
// Buffer DMA ADC circulaire (uint16). Half/Full = 1024 echantillons.

#define FFT_SIZE          4096
// FFT reelle embarquee. Resolution: FS/FFT_SIZE = 100000/4096 = 24.41 Hz/bin.

#define FS                100000.0f
// Frequence d'echantillonnage ADC (runtime via TIM2 ARR=1079).

#define DOPPLER_SCALE     161.0f
// Conversion fd->v pour radar 24.125 GHz.

#define MIN_BIN           2
#define MAX_BIN           250
// Fenetre de recherche pic: ~49 Hz a ~6.1 kHz.

#define SNR_THRESHOLD     15.0f
#define ADC_PP_THRESHOLD  100
// Double seuil de detection: relatif (SNR) + absolu (amplitude ADC).

#define UDP_DST_PORT      5555
#define UDP_SRC_PORT      5555
#define UDP_MAX_PAYLOAD   1400
#define UDP_HOP_SAMPLES   (ADC_BUF_SIZE / 2)  // 1024
// Streaming ADC brut en hops UDP fragmentes.
\end{lstlisting}

\subsection{Constantes DAC}

\begin{lstlisting}[style=cstyle,caption={Constantes du g\'en\'erateur de balayage DAC}]
#define DAC_BUF_SIZE  1024
// Taille du buffer circulaire pour le DAC.
// Le DMA envoie ces 1024 echantillons au DAC en boucle.
// A chaque fin de trame FFT, le buffer est re-rempli
// avec la nouvelle frequence du balayage.

#define SWEEP_F_START 30.0f
// Frequence de debut du balayage en Hz.
// 30 Hz correspond a environ 0.19 m/s en Doppler.

#define SWEEP_F_END   6000.0f
// Frequence de fin du balayage en Hz.
// 6 kHz deborde volontairement de la bande utile (3 kHz)
// pour verifier que MAX_BIN limite bien la detection.

#define SWEEP_TIME_MS 15000
// Duree totale du balayage : 15 secondes.
// Assez lent pour voir clairement les frequences defiler
// sur le terminal UART.
\end{lstlisting}

\newpage

% ══════════════════════════════════════════
% 3. VARIABLES GLOBALES
% ══════════════════════════════════════════
\section{Variables globales}

\subsection{Buffers et flags ADC (production)}

\begin{lstlisting}[style=cstyle,caption={D\'eclarations des buffers ADC et flags DMA}]
volatile uint16_t adc_buf[ADC_BUF_SIZE];
// Buffer circulaire rempli par le DMA depuis l'ADC.
// "volatile" car modifie par le DMA (materiel) en arriere-plan,
// pas par le CPU. Sans volatile, le compilateur pourrait
// optimiser les lectures et lire des valeurs perimees.
// uint16_t car l'ADC est 12 bits (valeurs 0 a 4095).

volatile uint8_t  adc_half_ready = 0;
volatile uint8_t  adc_full_ready = 0;
// Flags positionnes par les callbacks d'interruption DMA.
// adc_half_ready = 1 quand la 1ere moitie (indices 0..1023)
//                    est prete a etre lue
// adc_full_ready = 1 quand la 2eme moitie (indices 1024..2047)
//                    est prete a etre lue
// La boucle principale les remet a 0 apres traitement.
\end{lstlisting}

\subsection{Buffers FFT (production)}

\begin{lstlisting}[style=cstyle,caption={Buffers et \'etat de la FFT}]
float32_t fft_in[FFT_SIZE];        // 4096 echantillons d'entree
// On y accumule 4 x 1024 echantillons avant de lancer la FFT.
// float32_t = float simple precision (32 bits, ~7 chiffres).
// Les valeurs ADC (0-4095 entiers) sont converties en float
// lors de la copie depuis adc_buf.

float32_t fft_out[FFT_SIZE];       // sortie brute de la FFT
// Contient des nombres complexes entrelaces :
// [re0, im0, re1, im1, re2, im2, ...]
// Donc fft_out[2*k] = partie reelle du bin k
//      fft_out[2*k+1] = partie imaginaire du bin k

float32_t fft_mag[FFT_SIZE / 2];   // 2048 modules (spectre)
// Module = sqrt(re^2 + im^2) de chaque bin.
// fft_mag[k] represente l'amplitude a la frequence k*FS/FFT_SIZE
// Seuls les bins MIN_BIN a MAX_BIN sont examines.

float32_t hann_window[FFT_SIZE];   // coefficients de Hann
// Pre-calcules au demarrage par init_hann_window().
// Multiplies element par element avec fft_in avant la FFT.

arm_rfft_fast_instance_f32 fft_instance;
// Structure opaque de CMSIS-DSP contenant les "twiddle factors"
// (facteurs de rotation) pre-calcules pour une FFT de 4096 pts.
// Initialisee une seule fois par arm_rfft_fast_init_f32().

uint16_t fft_fill = 0;     // compteur : 0 a 4096
// Indique combien d'echantillons ont ete accumules dans fft_in.
// Incremente de 1024 a chaque moitie de buffer DMA.
// Quand il atteint FFT_SIZE (4096), on lance process_fft().
\end{lstlisting}

\subsection{Debug ADC (production)}

\begin{lstlisting}[style=cstyle,caption={Variables de diagnostic ADC}]
volatile uint16_t adc_min = 4095;  // valeur min vue
volatile uint16_t adc_max = 0;     // valeur max vue
volatile uint32_t adc_sum = 0;     // somme cumulee
volatile uint32_t adc_count = 0;   // nombre d'echantillons
// Ces statistiques sont collectees pendant la copie des
// echantillons ADC vers fft_in. Elles couvrent une trame
// complete (4096 echantillons).
//
// L'amplitude crete-a-crete (pp) = adc_max - adc_min
// sert de seuil absolu pour la detection :
//   - Signal reel de l'AFE : pp >> 100 LSB
//   - Bruit sur entree flottante : pp ~ 50 LSB
//   - Entree a la masse : pp ~ 10 LSB
//
// Remises a zero apres chaque FFT dans process_fft().
\end{lstlisting}

\subsection{Variables DAC}

\begin{lstlisting}[style=cstyle,caption={Variables du g\'en\'erateur de signal DAC}]
uint16_t dac_buf[DAC_BUF_SIZE];
// Buffer circulaire envoye au DAC par DMA.
// Contient 1024 echantillons d'une sinusoide a la
// frequence courante du balayage. Re-rempli a chaque
// trame FFT par update_dac_sweep().

float sweep_phase = 0.0f;
// Accumulateur de phase de la sinusoide DAC.
// Persiste entre les appels a update_dac_sweep() pour
// garantir la CONTINUITE DE PHASE : quand on re-remplit
// le buffer, la sinusoide reprend exactement la ou elle
// s'etait arretee. Sans cela, il y aurait une
// discontinuite qui creerait des harmoniques parasites.

uint32_t sweep_start_tick = 0;
// Timestamp (en ms) du debut du balayage.
// HAL_GetTick() retourne le compteur SysTick (1 kHz).
// Le temps ecoule = HAL_GetTick() - sweep_start_tick
// determine la frequence courante dans le balayage log.
\end{lstlisting}

\newpage

% ══════════════════════════════════════════
% 4. FONCTIONS UTILITAIRES
% ══════════════════════════════════════════
\section{Fonctions utilitaires}

\subsection{init\_hann\_window() -- Production}

\begin{lstlisting}[style=cstyle,caption={Pr\'e-calcul de la fen\^etre de Hann}]
void init_hann_window(void)
{
  for (int i = 0; i < FFT_SIZE; i++)
    hann_window[i] = 0.5f * (1.0f - arm_cos_f32(
                       2.0f * PI * i / (FFT_SIZE - 1)));
}
// Formule : w[i] = 0.5 * (1 - cos(2*pi*i / (N-1)))
//
// POURQUOI LE FENETRAGE ?
// ========================
// La FFT suppose que le signal est periodique sur la
// fenetre d'analyse. Or notre bloc de 4096 echantillons
// est une "coupe" arbitraire du signal continu.
// Les bords du bloc ont generalement des discontinuites
// qui creent des "fuites spectrales" : l'energie d'un
// pic se repand sur les bins voisins.
//
// La fenetre de Hann attenue progressivement les bords
// vers zero :
//   - w[0] = w[N-1] = 0  (bords a zero)
//   - w[N/2] = 1.0        (centre inchange)
//
// Resultat : les fuites spectrales sont reduites d'un
// facteur ~100 (de -13 dB a -32 dB pour le premier
// lobe lateral). Le pic est un peu plus large (2 bins
// au lieu de 1) mais beaucoup plus propre.
//
// arm_cos_f32() utilise une LUT + interpolation,
// plus rapide que cosf() de math.h sur Cortex-M7.
\end{lstlisting}

\subsection{get\_sweep\_freq() -- DAC}

\begin{lstlisting}[style=cstyle,caption={Calcul de la fr\'equence instantan\'ee du balayage}]
float get_sweep_freq(void)
{
  uint32_t elapsed = HAL_GetTick() - sweep_start_tick;
  // Temps ecoule en ms depuis le debut du balayage

  if (elapsed >= SWEEP_TIME_MS)
  {
    sweep_start_tick = HAL_GetTick();  // restart
    sweep_phase = 0.0f;
    elapsed = 0;
    // Le balayage a termine ses 15 secondes : on
    // recommence depuis le debut (30 Hz).
    // La phase est remise a zero pour repartir proprement.
  }

  float t_ratio = (float)elapsed / (float)SWEEP_TIME_MS;
  // t_ratio varie lineairement de 0.0 a 1.0

  return SWEEP_F_START * powf(SWEEP_F_END / SWEEP_F_START,
                              t_ratio);
  // BALAYAGE LOGARITHMIQUE :
  // freq(t) = F_start * (F_end / F_start)^(t/T)
  //
  // A t=0     : freq = 30 * (200)^0     = 30 Hz
  // A t=T/2   : freq = 30 * (200)^0.5   = 424 Hz
  // A t=T     : freq = 30 * (200)^1     = 6000 Hz
  //
  // Pourquoi logarithmique ? Chaque octave (doublement de
  // frequence) recoit le meme temps de balayage.
  // Lineaire donnerait 90% du temps entre 600 et 6000 Hz
  // et seulement 10% entre 30 et 600 Hz.
}
\end{lstlisting}

\subsection{update\_dac\_sweep() -- DAC}

\begin{lstlisting}[style=cstyle,caption={Remplissage du buffer DAC avec la sinuso\"ide}]
void update_dac_sweep(void)
{
  float freq = get_sweep_freq();
  // Frequence actuelle du balayage (30 a 6000 Hz)

  float phase_inc = 2.0f * PI * freq / FS;
  // Increment de phase par echantillon.
  // A 100 kHz d'echantillonnage et 500 Hz de signal :
  //   phase_inc = 2*pi*500/100000 = 0.0314 rad
  //   soit 100000/500 = 200 echantillons par periode

  for (int i = 0; i < DAC_BUF_SIZE; i++)
  {
    dac_buf[i] = (uint16_t)(2048.0f + 800.0f
                            * sinf(sweep_phase));
    // DAC 12 bits : plage 0 a 4095
    // 2048 = milieu (1.65V sur 3.3V)
    // +/- 800 LSB = +/- 0.64V d'amplitude
    // Signal total : 1.01V a 2.29V

    sweep_phase += phase_inc;
    // L'accumulateur de phase avance d'un pas.
    // Cle de la CONTINUITE DE PHASE : la phase
    // n'est JAMAIS remise a zero entre deux buffers,
    // donc la sinusoide est parfaitement continue.
  }

  if (sweep_phase > 2.0f * PI * 1000.0f)
    sweep_phase = fmodf(sweep_phase, 2.0f * PI);
  // Prevention du debordement float : quand la phase
  // accumule depasse 6283 rad (~1000 tours), on la
  // replie dans [0, 2*pi]. sinf() perd en precision
  // pour les tres grandes valeurs de phase.
}
\end{lstlisting}

\newpage

% ══════════════════════════════════════════
% 5. PROCESS_FFT
% ══════════════════════════════════════════
\section{Fonction process\_fft() -- Le coeur du traitement}

Cette fonction est appel\'ee chaque fois que 4096 \'echantillons ont \'et\'e accumul\'es. Elle constitue le c\oe ur du traitement du signal.

\subsection{\'Etape 1 : Suppression DC}

\begin{lstlisting}[style=cstyle,caption={Suppression de la composante continue}]
void process_fft(void)
{
  float32_t mean;
  arm_mean_f32(fft_in, FFT_SIZE, &mean);
  // Calcule la moyenne arithmetique des 4096 echantillons.
  // Pour un ADC centre a 1.65V : mean ~ 2048 LSB

  for (int i = 0; i < FFT_SIZE; i++)
    fft_in[i] -= mean;
  // Soustrait la moyenne de chaque echantillon.
  // Resultat : le signal oscille autour de 0 au lieu de 2048.
  //
  // POURQUOI ?
  // La composante DC (offset constant) se traduit en un pic
  // enorme au bin 0 de la FFT. Ce pic "deborde" sur les
  // bins voisins (leakage) et masque les petits signaux
  // Doppler. En soustrayant la moyenne, on supprime ce pic.
\end{lstlisting}

\subsection{\'Etape 2 : Fen\^etrage de Hann}

\begin{lstlisting}[style=cstyle,caption={Application de la fen\^etre de Hann}]
  arm_mult_f32(fft_in, hann_window, fft_in, FFT_SIZE);
  // Multiplication element par element :
  //   fft_in[i] = fft_in[i] * hann_window[i]
  //
  // arm_mult_f32 est optimisee CMSIS-DSP : elle traite
  // 4 multiplications par cycle grace au pipeline FPU
  // du Cortex-M7. Pour 4096 elements, le cout reste
  // tres faible devant la FFT elle-meme.
  // cycles soit ~2.4 us a 216 MHz.
  //
  // Le resultat est un signal dont les bords valent 0,
  // pret pour la FFT sans fuites spectrales excessives.
\end{lstlisting}

\subsection{\'Etape 3 : FFT et magnitude}

\begin{lstlisting}[style=cstyle,caption={FFT rapide r\'eelle et calcul du spectre d'amplitude}]
  arm_rfft_fast_f32(&fft_instance, fft_in, fft_out, 0);
  // Calcule la FFT reelle 4096 points.
  //
  // "Reelle" signifie que l'entree est purement reelle
  // (pas de partie imaginaire), ce qui permet d'utiliser
  // un algorithme 2x plus rapide que la FFT complexe.
  //
  // Le parametre 0 = FFT directe (1 = FFT inverse).
  //
  // Sortie : 4096 floats = 2048 nombres complexes
  //   fft_out[0] = DC (partie reelle)
  //   fft_out[1] = Nyquist (partie reelle)
  //   fft_out[2k], fft_out[2k+1] = re, im du bin k
  //
  // Complexite : O(N*log2(N)) = 4096*12 = 49152 operations
  // Duree typique : ~quelques ms sur Cortex-M7 @216 MHz (selon charge systeme)

  arm_cmplx_mag_f32(fft_out, fft_mag, FFT_SIZE / 2);
  // Pour chaque paire (re, im), calcule :
  //   mag = sqrt(re*re + im*im)
  //
  // Le resultat fft_mag[k] est l'AMPLITUDE au bin k.
  // La frequence correspondante est :
  //   f = k * FS / FFT_SIZE = k * 24.414 Hz
  //
  // Exemple : fft_mag[20] = amplitude a 488.3 Hz
\end{lstlisting}

\subsection{\'Etape 4 : Recherche du pic}

\begin{lstlisting}[style=cstyle,caption={Recherche du maximum dans la fen\^etre de fr\'equence}]
  float32_t peak_val;
  uint32_t peak_idx;
  arm_max_f32(&fft_mag[MIN_BIN], MAX_BIN - MIN_BIN,
              &peak_val, &peak_idx);
  peak_idx += MIN_BIN;
  // Cherche la valeur maximale dans fft_mag[2..249].
  // arm_max_f32 retourne la valeur ET l'indice du max.
  // On ajoute MIN_BIN car la recherche commence a l'offset
  // MIN_BIN, pas a l'indice 0 du tableau.
  //
  // peak_idx * 24.414 = frequence Doppler detectee
  // peak_val = amplitude du pic (en unites lineaires)
\end{lstlisting}

\subsection{\'Etape 5 : Estimation du bruit (noise floor)}

\begin{lstlisting}[style=cstyle,caption={Calcul du plancher de bruit en excluant le pic}]
  float32_t noise_sum = 0;
  int noise_count = 0;
  for (int i = MIN_BIN; i < MAX_BIN; i++)
  {
    if (i < (int)peak_idx - 5 || i > (int)peak_idx + 5)
    {
      noise_sum += fft_mag[i];
      noise_count++;
    }
  }
  float32_t noise_avg = (noise_count > 0) ?
                         noise_sum / noise_count : 1e-12f;
  // On calcule la moyenne de TOUS les bins de la fenetre
  // SAUF les 11 bins autour du pic (pic +/- 5 bins).
  //
  // POURQUOI EXCLURE LE PIC ?
  // Si le signal est fort, inclure le pic et ses lobes
  // lateraux dans la moyenne du bruit gonflerait
  // artificiellement le "bruit", reduisant le SNR calcule.
  // En excluant +/- 5 bins, on mesure le vrai bruit de fond.
  //
  // 1e-12f evite une division par zero si noise_count = 0.
\end{lstlisting}

\subsection{\'Etape 6 : Calcul du SNR, d\'ecision et sorties}

\begin{lstlisting}[style=cstyle,caption={D\'ecision de d\'etection + sorties UART/OLED/LED}]
  float32_t peak_db = 20.0f * log10f(peak_val + 1e-12f);
  float32_t noise_db = 20.0f * log10f(noise_avg + 1e-12f);
  float32_t snr = peak_db - noise_db;

  float32_t freq = (float32_t)peak_idx * FS / FFT_SIZE;
  float32_t speed_ms = freq / DOPPLER_SCALE;
  float32_t speed_kmh = speed_ms * 3.6f;

  uint16_t adc_pp = adc_max - adc_min;
  adc_min = 4095; adc_max = 0; adc_sum = 0; adc_count = 0;

  if (snr > SNR_THRESHOLD && adc_pp > ADC_PP_THRESHOLD)
  {
    HAL_GPIO_WritePin(GPIOB, LD1_Pin, GPIO_PIN_SET);
    HAL_GPIO_WritePin(GPIOB, LD3_Pin, GPIO_PIN_RESET);
    last_speed_kmh = speed_kmh;
    last_speed_ms  = speed_ms;
    last_freq      = freq;
    last_snr       = snr;
    last_detected  = 1;
  }
  else
  {
    HAL_GPIO_WritePin(GPIOB, LD1_Pin, GPIO_PIN_RESET);
    HAL_GPIO_WritePin(GPIOB, LD3_Pin, GPIO_PIN_SET);
    last_snr      = snr;
    last_detected = 0;
  }

  // UART binaire legacy:
  // 0xAA55 + n_samples + stream_buf[1024]
  // 0xBB66 + freq + speed_kmh + speed_ms + snr + det
\end{lstlisting}

\newpage

% ══════════════════════════════════════════
% 6. MAIN - INITIALISATION
% ══════════════════════════════════════════
\section{Fonction main() -- S\'equence d'initialisation}

\subsection{Cache, HAL et horloge}

\begin{lstlisting}[style=cstyle,caption={D\'emarrage du syst\`eme}]
int main(void)
{
  // === CACHE ===
  SCB_EnableICache();
  // Active le cache d'instructions (I-Cache) du Cortex-M7.
  // Accelere l'execution en evitant d'aller chercher chaque
  // instruction en Flash (latence ~7 cycles a 216 MHz).
  //
  // Le D-Cache (donnees) est VOLONTAIREMENT DESACTIVE :
  // le DMA ecrit directement en RAM sans passer par le cache.
  // Si le D-Cache etait actif, le CPU lirait des donnees
  // perimees (cachees) au lieu des donnees DMA fraiches.
  // Solution alternative : invalidation manuelle du cache
  // apres chaque transfert DMA, mais c'est plus complexe.

  HAL_Init();
  // Initialise le HAL, configure SysTick a 1 kHz (1 tick/ms)
  // pour HAL_Delay() et HAL_GetTick().

  SystemClock_Config();
  // Configure la PLL : HSE 8 MHz -> SYSCLK 216 MHz
  // PLL : M=4, N=216, P=2 -> 8/4*216/2 = 216 MHz
  // APB1 = 216/4 = 54 MHz (timers APB1 = 108 MHz)
  // APB2 = 216/2 = 108 MHz (timers APB2 = 216 MHz)
  // Note : les timers sur APB1/APB2 tournent a 2x
  // la frequence du bus si le prescaler bus > 1.
\end{lstlisting}

\subsection{Initialisation des p\'eriph\'eriques et d\'emarrage}

\begin{lstlisting}[style=cstyle,caption={Init des p\'eriph\'eriques et d\'emarrage de l'acquisition}]
  // === PERIPHERIQUES (genere par CubeMX) ===
  MX_GPIO_Init();           // LEDs (PB0, PB7, PB14), bouton
  MX_DMA_Init();            // DMA2 Stream0 pour ADC
  MX_ETH_Init();            // Ethernet utilise pour streaming UDP ADC
  MX_USART3_UART_Init();    // UART3 a 115200 (modifie apres)
  MX_USB_OTG_FS_PCD_Init(); // USB (non utilise)
  MX_ADC1_Init();           // ADC1 CH0, TIM2 trigger, 12 bit
  MX_DAC_Init();            // DAC CH1, TIM4 trigger
  MX_TIM2_Init();           // TIM2 CubeMX, puis reconfig runtime a 100 kHz
  MX_TIM4_Init();           // TIM4 : init CubeMX (modifie apres)

  // === UART RAPIDE ===
  huart3.Init.BaudRate = 921600;
  if (HAL_UART_Init(&huart3) != HAL_OK) Error_Handler();
  // CubeMX initialise UART a 115200 bauds.
  // 115200 est trop lent : une trame FFT toutes les 41 ms
  // avec ~60 caracteres = ~1460 char/s.
  // A 115200 bauds (11520 char/s) ca passerait, mais sans
  // marge pour du debug plus verbeux.
  // 921600 = 8x plus rapide, aucun souci de debit.

  // === FFT ===
  arm_rfft_fast_init_f32(&fft_instance, FFT_SIZE);
  // Pre-calcule les "twiddle factors" pour la FFT 4096 pts.
  // Ce sont des valeurs cos/sin utilisees a chaque etage
  // de l'algorithme papillon (butterfly) de la FFT.
  // Cette init prend ~0.1 ms et n'est faite qu'une fois.

  init_hann_window();
  // Pre-calcule les 4096 coefficients de la fenetre de Hann.

  // === DEMARRAGE ADC + DMA + TIMER ===
  HAL_ADC_Start_DMA(&hadc1, (uint32_t *)adc_buf, ADC_BUF_SIZE);
  // Lance l'ADC en mode DMA circulaire :
  // - Le DMA ecrira dans adc_buf[0..1023] en boucle
  // - Interruption a mi-parcours (1024) et a la fin (2048)
  // - Le cast (uint32_t*) est requis par le HAL meme si
  //   les donnees sont 16 bits (le DMA est configure en
  //   half-word dans le MSP)

  HAL_TIM_Base_Start(&htim2);
  // Demarre TIM2. A partir de cet instant, l'ADC
  // echantillonne a 100 kHz. Les donnees arrivent dans
  // adc_buf via DMA automatiquement.
\end{lstlisting}

\subsection{Configuration du DAC}

\begin{lstlisting}[style=cstyle,caption={Init manuelle du DMA DAC et d\'emarrage du balayage}]
  // === DESACTIVATION IRQ TIM4 ===
  HAL_NVIC_DisableIRQ(TIM4_IRQn);
  // TIM4 etait utilise pour un toggle LED a 60 Hz.
  // On le reutilise pour piloter le DAC a 100 kHz.
  // Si l'IRQ restait active, le CPU recevrait 100000
  // interruptions/s rien que pour TIM4, ce qui
  // saturerait completement le processeur.

  sweep_start_tick = HAL_GetTick();
  update_dac_sweep();  // remplir le buffer initial

  // === CONFIGURATION MANUELLE DMA1 STREAM5 CH7 ===
  // Le .ioc ne configure pas le DMA pour le DAC,
  // donc on le fait manuellement dans le code utilisateur.
  __HAL_RCC_DMA1_CLK_ENABLE();            // horloge DMA1
  hdma_dac1.Instance = DMA1_Stream5;       // stream 5
  hdma_dac1.Init.Channel = DMA_CHANNEL_7;  // canal 7
  // Le mapping stream/channel est fixe par le hardware STM32 :
  // DMA1_Stream5_Channel7 = DAC1 (voir Reference Manual table)
  hdma_dac1.Init.Direction = DMA_MEMORY_TO_PERIPH;
  // Sens : memoire (dac_buf) -> peripherique (registre DAC)
  hdma_dac1.Init.PeriphInc = DMA_PINC_DISABLE;
  // L'adresse peripherique (DAC->DHR12R1) ne change pas
  hdma_dac1.Init.MemInc = DMA_MINC_ENABLE;
  // L'adresse memoire avance de +2 octets a chaque transfert
  hdma_dac1.Init.PeriphDataAlignment = DMA_PDATAALIGN_HALFWORD;
  hdma_dac1.Init.MemDataAlignment = DMA_MDATAALIGN_HALFWORD;
  // 16 bits (half-word) car le DAC est 12 bits aligne a droite
  hdma_dac1.Init.Mode = DMA_CIRCULAR;
  // Mode circulaire : quand il arrive a la fin de dac_buf,
  // le DMA recommence au debut automatiquement.
  hdma_dac1.Init.Priority = DMA_PRIORITY_LOW;
  hdma_dac1.Init.FIFOMode = DMA_FIFOMODE_DISABLE;
  HAL_DMA_Init(&hdma_dac1);
  __HAL_LINKDMA(&hdac, DMA_Handle1, hdma_dac1);
  // Lie le handle DMA au handle DAC (requis par le HAL)

  // === TIMER 4 A 100 kHz ===
  __HAL_TIM_SET_PRESCALER(&htim4, 0);
  __HAL_TIM_SET_AUTORELOAD(&htim4, 1079);
  // 108 MHz / (0+1) / (1079+1) = 100000 Hz
  // Meme frequence que l'ADC pour que le DAC genere
  // exactement FS echantillons par seconde.
  HAL_TIM_Base_Start(&htim4);

  HAL_DAC_Start_DMA(&hdac, DAC_CHANNEL_1,
                    (uint32_t *)dac_buf,
                    DAC_BUF_SIZE, DAC_ALIGN_12B_R);
  // Lance le DAC en mode DMA circulaire.
  // Chaque tick TIM4 (100 kHz), le DMA envoie un
  // echantillon de dac_buf au DAC.
  // Le signal sort physiquement sur la pin PA4.
\end{lstlisting}

\newpage

% ══════════════════════════════════════════
% 7. BOUCLE PRINCIPALE
% ══════════════════════════════════════════
\section{Boucle principale (while)}

\begin{lstlisting}[style=cstyle,caption={Pattern double-buffer dans la boucle principale}]
  while (1)
  {
    // === PREMIERE MOITIE DU BUFFER (indices 0..1023) ===
    if (adc_half_ready)
    {
      adc_half_ready = 0;  // acquitter le flag

      for (int i = 0; i < ADC_BUF_SIZE / 2; i++)
      {
        uint16_t val = adc_buf[i];   // lire un echantillon
        fft_in[fft_fill + i] = (float32_t)val;  // copier

        // Collecter les stats ADC pendant la copie
        if (val < adc_min) adc_min = val;
        if (val > adc_max) adc_max = val;
        adc_sum += val;
        adc_count++;
      }

      fft_fill += ADC_BUF_SIZE / 2;  // +1024

      if (eth_link_up)
        udp_send_adc_frame(udp_hop_buf, UDP_HOP_SAMPLES, udp_seq_counter++);

      if (fft_fill >= FFT_SIZE)  // 4096 atteint ?
      {
        fft_fill = 0;       // reset pour la prochaine trame
        process_fft();       // lancer le traitement FFT
        update_dac_sweep();  // rafraichir le DAC
        HAL_GPIO_TogglePin(GPIOB, LD2_Pin);  // LED bleue
      }
    }

    // === SECONDE MOITIE DU BUFFER (indices 1024..2047) ===
    if (adc_full_ready)
    {
      adc_full_ready = 0;  // acquitter le flag

      for (int i = 0; i < ADC_BUF_SIZE / 2; i++)
      {
        uint16_t val = adc_buf[ADC_BUF_SIZE / 2 + i];  // +1024
        fft_in[fft_fill + i] = (float32_t)val;

        if (val < adc_min) adc_min = val;
        if (val > adc_max) adc_max = val;
        adc_sum += val;
        adc_count++;
      }

      fft_fill += ADC_BUF_SIZE / 2;

      if (eth_link_up)
        udp_send_adc_frame(udp_hop_buf, UDP_HOP_SAMPLES, udp_seq_counter++);

      if (fft_fill >= FFT_SIZE)
      {
        fft_fill = 0;
        process_fft();
        update_dac_sweep();  // rafraichir le DAC
        HAL_GPIO_TogglePin(GPIOB, LD2_Pin);
      }
    }

    // OLED a 5 Hz
    if (HAL_GetTick() - oled_last_update >= OLED_UPDATE_MS)
      SSD1306_ShowRadar(last_speed_kmh, last_speed_ms,
                        last_freq, last_snr, last_detected);

    // Verification lien Ethernet a 1 Hz
    if (HAL_GetTick() - eth_check_tick >= 1000) { ... }
  }
}
\end{lstlisting}

\subsubsection*{Explication du pattern double-buffer}

Le DMA remplit \code{adc\_buf[2048]} en continu et de mani\`ere circulaire. Le CPU et le DMA ne doivent jamais acc\'eder \`a la m\^eme moiti\'e en m\^eme temps :

\begin{center}
\small
\begin{tabular}{c c c}
\toprule
\textbf{Phase} & \textbf{DMA \'ecrit dans} & \textbf{CPU lit} \\
\midrule
Half-transfer & indices 1024..2047 & indices 0..1023 \\
Full-transfer & indices 0..1023    & indices 1024..2047 \\
\bottomrule
\end{tabular}
\end{center}

Cadences principales :
\begin{itemize}[leftmargin=2em]
  \item Hop DMA/UDP : \(1024 / 100000 = 10.24\,ms\) (\(\sim 97.6\) hops/s)
  \item Trame FFT : \(4096 / 100000 = 40.96\,ms\) (\(\sim 24.4\) FFT/s)
  \item OLED : 5 Hz ; v\'erification lien Ethernet : 1 Hz
\end{itemize}

\newpage

% ══════════════════════════════════════════
% 8. CALLBACKS DMA
% ══════════════════════════════════════════
\section{Callbacks DMA}

\begin{lstlisting}[style=cstyle,caption={Callbacks d'interruption DMA pour l'ADC}]
void HAL_ADC_ConvHalfCpltCallback(ADC_HandleTypeDef *hadc)
{
  if (hadc->Instance == ADC1) adc_half_ready = 1;
}
// Appelee par le HAL depuis l'ISR DMA2_Stream0.
// Declenchee quand le DMA a rempli adc_buf[0..1023].
// Ne fait QU'une seule chose : lever un flag.
//
// REGLE D'OR DES ISR :
// Les routines d'interruption doivent etre les plus courtes
// possible. Le traitement lourd (FFT, UART) se fait dans
// la boucle principale, pas dans l'ISR.

void HAL_ADC_ConvCpltCallback(ADC_HandleTypeDef *hadc)
{
  if (hadc->Instance == ADC1) adc_full_ready = 1;
}
// Idem, mais declenchee quand adc_buf[1024..2047] est rempli.
// Le DMA recommence ensuite a l'indice 0 (mode circulaire).

/* TIM4 callback desactive :
 * TIM4 pilote maintenant le DAC DMA a 100 kHz.
 * L'ancien callback (toggle LED a 60 Hz) a ete retire.
 * Si l'IRQ TIM4 etait active, 100000 interruptions/s
 * consommeraient ~30% du CPU pour rien.
 */
\end{lstlisting}

\newpage

% ══════════════════════════════════════════
% 9. CONFIG CubeMX
% ══════════════════════════════════════════
\section{Configuration des p\'eriph\'eriques (CubeMX)}

\subsection{ADC1 -- MX\_ADC1\_Init()}

\begin{lstlisting}[style=cstyle,caption={Configuration de l'ADC1}]
  hadc1.Instance = ADC1;
  hadc1.Init.ClockPrescaler = ADC_CLOCK_SYNC_PCLK_DIV4;
  // Horloge ADC = APB2 / 4 = 108 / 4 = 27 MHz
  // Frequence max de l'ADC STM32F7 = 36 MHz, on est ok.

  hadc1.Init.Resolution = ADC_RESOLUTION_12B;
  // 12 bits : valeurs de 0 a 4095

  hadc1.Init.ScanConvMode = ADC_SCAN_DISABLE;
  hadc1.Init.ContinuousConvMode = DISABLE;
  // Un seul canal, pas de scan. Le mode continu est
  // desactive car c'est TIM2 qui declenche chaque conversion.

  hadc1.Init.ExternalTrigConvEdge = ADC_EXTERNALTRIGCONVEDGE_RISING;
  hadc1.Init.ExternalTrigConv = ADC_EXTERNALTRIGCONV_T2_TRGO;
  // Declenchement sur front montant du TRGO de TIM2.
  // A chaque evenement Update de TIM2 (runtime: 100 kHz), une
  // conversion ADC demarre automatiquement.

  hadc1.Init.DMAContinuousRequests = ENABLE;
  // Le DMA continue de transferer apres chaque conversion
  // (mode circulaire). Sans ca, le DMA s'arreterait apres
  // le premier buffer.

  // Canal echantillonne :
  sConfig.Channel = ADC_CHANNEL_0;  // PA0
  sConfig.SamplingTime = ADC_SAMPLETIME_15CYCLES;
  // Temps de conversion total =
  //   (15 cycles sampling + 12 cycles conversion) / 27 MHz
  //   = 27 / 27e6 = 1.0 us (ordre de grandeur)
  // A 100 kHz (10 us entre chaque conversion), largement ok.
\end{lstlisting}

\subsection{TIM2 -- Horloge ADC (reconfigur\'ee \`a 100\,kHz)}

\begin{lstlisting}[style=cstyle,caption={Configuration du timer d'\'echantillonnage}]
  htim2.Instance = TIM2;
  htim2.Init.Prescaler = 0;
  htim2.Init.Period = 2159;  // ARR = 2159
  // Valeur CubeMX (50 kHz). Dans USER CODE, ARR est force a 1079
  // pour obtenir 100 kHz runtime :
  //   f = 108e6 / (PSC+1) / (ARR+1) = 108e6 / 1 / 1080 = 100 kHz.

  sMasterConfig.MasterOutputTrigger = TIM_TRGO_UPDATE;
  // L'evenement Update est route vers la sortie TRGO.
  // Le TRGO de TIM2 est connecte a l'entree trigger
  // de l'ADC1 (T2_TRGO), ce qui declenche une conversion.
\end{lstlisting}

\newpage

% ══════════════════════════════════════════
% 10. RESUME FINAL
% ══════════════════════════════════════════
\section{R\'esum\'e final}

Le firmware actuel combine :
\begin{itemize}[leftmargin=2em]
  \item acquisition ADC \`a 100~kHz + DMA circulaire,
  \item traitement FFT 4096 points avec d\'etection Doppler,
  \item streaming UDP hops vers le PC (port 5555),
  \item t\'el\'em\'etrie UART legacy et affichage OLED,
  \item g\'en\'eration DAC synchronis\'ee (TIM4 + DMA) conserv\'ee dans l'architecture.
\end{itemize}

\vfill
\begin{center}
\textcolor{darkgray}{\small\textit{Fin du document -- Projet L3 EEA Radar Doppler 24\,GHz}}
\end{center}

\end{document}
